\scanpage{1}
% not manuscript: couverture institutionnelle de l'Université de Montpellier et de l'IMAG.

\scanpage{2}
\unsure{Pages|Page} 650--787 + table of contents

\uline{N.B.} Page 650 appears twice (numbering error)

\margin{1}

% Pied de page imprimé de l'archive.
\begin{center}
\small
1\\
\texttt{https://grothendieck.umontpellier.fr}\\
Université de Montpellier, UMR CNRS 5149, Case courrier 051, Place Eugène Bataillon\\
34095 Montpellier cedex 5 - France
\end{center}

\scanpage{3}
50. Schematic envelopes of discrete groups: metabelian case \hfill 650

1°) Unipotent sections \hfill same

2°) Attempt at a \unsure{characterization} of quasi-unipotent sections \hfill $652'$

3°) Group scheme associated with certain extensions of discrete groups (abelian kernel) \hfill $\unsure{655'}$

4°) Schematic envelope of a metabelian group \hfill $658'$

5°) Universal characterizatio\textsuperscript{n} of the schematic envelope \hfill $\unsure{662'}$

51. Example: the Fermat tower. \hfill 664

1°) The groups $\mathfrak S=\mathrm{Sl}(2,\mathbb Z)/[\Pi_0,\Pi_0]$, and $\underline{\mathfrak S}$ \hfill same

2°) $\unsure{\underline Z|Z}_{\rho,\sigma}$: provisional version \hfill $665'$

3°) The groups $\Phi_N^*$ and variants \hfill 673

4°) The sequence
\diagram{\begin{tikzcd}[column sep=small]
1 \arrow[r] & \underline N^{0+}_{\rho,\sigma} \arrow[r] & \underline N_{\rho,\sigma} \arrow[r] & \Phi_{12}^{*} \arrow[r] & 1
\end{tikzcd}},
and relations with $\mathrm{Sl}(2,\mathbb Z/12\mathbb Z)$ \hfill 677

5°) $\underline M_{\rho,\sigma}$, $\underline M^{\unsure{\natural|\sharp}}_{\rho,\sigma}$ via $N_{\rho,\sigma}$ \hfill 682

6°) \deleted{$\unsure{\underline Z_{\rho,\sigma}}$} \unsure{sequences}
\diagram{\begin{tikzcd}[column sep=small]
1 \arrow[r] & \mathbb G_a \arrow[r] & \underline Z_{\rho,\sigma} \arrow[r] & \widehat\Phi_{12} \arrow[r] & 1
\end{tikzcd}},
and
\diagram{\begin{tikzcd}[column sep=small]
1 \arrow[r] & \underline{\mathfrak S}^{+} \arrow[r] & \underline G_{\rho,\sigma} \arrow[r] & \widehat\Phi_{12} \arrow[r] & 1
\end{tikzcd}}
\hfill 684

7°) $\underline N_\rho$, $\underline N_\rho^{\unsure{\natural|\sharp}}$ and so on \hfill 686

8°) $\underline N_\sigma$, $\underline N_\sigma^{\unsure{\natural|\sharp}}$ and so on \hfill 689

9°) Return to $\underline M_{\rho,\sigma}$, $\underline M_{\rho,\sigma}^{\unsure{\natural|\sharp}}$ and so on \hfill 690

10°) $\underline M_{\rho,\sigma}\unsure{(0)|[0]}$, $\underline M_{\rho,\sigma}\unsure{[\tfrac12]|(\tfrac12)}$, $\underline M_{\rho,\sigma}(-\bar j)$ and so on \hfill $691'$

11°) The homomorphism $M\longrightarrow M_{\rho,\sigma}=\underline M_{\rho,\sigma}(\widehat{\mathbb Z})^{\wedge}$, and relations with the tower of Fermat curves: statement of the problem. \hfill 695

12°) Splittings of the extension $G_{\rho,\sigma}/\mu_G$ of $\widehat{\mathbb Z}^{*}\times\mathfrak S_3$ by $\widehat M$. \hfill 700

13°) \deleted{\unsure{Digression on}}\inserted{Characterization of} the Fermat tower; \deleted{and} \deleted{\unsure{conjectures}}\inserted{relations} of compatibility, with $M\longrightarrow G_{\rho,\sigma}$ \hfill \deleted{\unsure{703}}\inserted{707}

14°) Relation with generalized Jacobians \hfill $\unsure{720'}$

15°) Interpretation of $\unsure{\Gamma_{\mathbb Q}}\longrightarrow N_\sigma^{\unsure{\natural|\sharp}}\simeq\unsure{\widehat{\mathbb Z}}^{*}\mathbin{.}\widehat{\mathbb Z}$, and digression on \unsure{Kummer} extensions \hfill $725'$

16°) Reflections on Fermat's theorem \hfill 733

\margin{2}

% Pied de page imprimé de l'archive.
\begin{center}
\small
2\\
\texttt{https://grothendieck.umontpellier.fr}\\
Université de Montpellier, UMR CNRS 5149, Case courrier 051, Place Eugène Bataillon\\
34095 Montpellier cedex 5 - France
\end{center}

\scanpage{4}
52. \deleted{Schematic envelopes of L-groups and $\Pi L$-groups.}\inserted{Lie approximations of Galois-Teichmüller systems.}

1°) Notion of L- and $\Pi L$-group: discrete case \hfill 737

2°) Profinite case \hfill 744

3°) \deleted{\unsure{Discretification}, and \unsure{admissible homomorphisms},}\inserted{Lie approximations of systems} of Galois-Teichmüller \hfill 748

4°) Lie approximation induced by a ``geometric'' Lie approximation \hfill 758

\uline{53}. The special Teichmüller group $\unsure{S\mathcal T}_{1,1}\simeq\mathrm{Gl}(2,\mathbb Z)^{\sim}$, and the group $SN_{1,1}$ \hfill 760

1°) Differential action and special Teichmüller group (provisional statement) \hfill same

2°) The groups $\mathrm{Gl}(2,\mathbb R)^{\pm1\,\sim}$, $\mathrm{Sl}(2,\mathbb R)^{\sim}$ \hfill $761'$

3°) The group $\mathrm{Gl}(2,\mathbb Z)^{\sim}$ \hfill 763

4°) The isomorphism $\mathrm{Gl}(2,\mathbb Z)^{\sim}\simeq\unsure{S\mathcal T}_{1,1}$ \hfill $\unsure{766'|765'}$

5°) Digression on coverings of germs of spaces, and relations between topological germs and differentiable germs of surfaces \hfill $768'$

6°)

% Right-hand page of the spread: blank.

\margin{3}

% Pied de page imprimé de l'archive.
\begin{center}
\small
3\\
\texttt{https://grothendieck.umontpellier.fr}\\
Université de Montpellier, UMR CNRS 5149, Case courrier 051, Place Eugène Bataillon\\
34095 Montpellier cedex 5 - France
\end{center}

\scanpage{5}
\mspage{650\ \textcircled{bis}}

50. \uline{Schematic envelopes of discrete groups.}

\deleted{\unsure{2}}(\uline{Metabelian case})

1.\deleted{\unsure{g}} Let $G$ be a group scheme locally of finite presentation
over a base scheme $S$ (the case \unsure{important} for us will be
$S=\operatorname{Spec}(\mathbb Z)$). A section $g$ of $G$ over $S$ is said to be
\uline{unipotent} if on every geometric fiber,
\deleted{\unsure{of}} $G_{\bar s}$, the element
$g_{\bar s}\in G_{\bar s}(\unsure{\bar k|k^a})$
$(k=k(\bar s))$ is unipotent, i.e. if there exists a group homomorphism
\margin{\uline{Unipotent \unsure{sections}}}
\[
(1)\qquad \mathbb G_{ak}\xrightarrow{\varphi_{\bar s}}G_k
\qquad k=k(\bar s)
\]
such that
\[
(2)\qquad g_{\bar s}=\varphi_{\bar s}(1)
\]

We say that $g$ is quasi-unipotent if for every $\bar s$, there exists
$n\in\mathbb N^*$ such that $g_{\bar s}^{n}$ is unipotent.
If the characteristic of $\bar s$ is $p>0$, then this also \uline{means} that
$g_{\bar s}$ is of \uline{finite order}. We see that if $S$ is
quasi-compact, then $g$ is quasi-unipotent iff
$\exists n\in\mathbb N^*$ such that $g^n$ is unipotent.

\deleted{\uline{Theorem} -- Let $G$ \inserted{of groups}
be locally of finite presentation over quasi-compact $S$,
$\unsure{g|G}$ \unsure{if|which} is \unsure{unipotent|unipotent}, \unsure{for a}}

\deleted{\unsure{?}} For $g$ to be unipotent, it \uline{suffices} that there exist a homomorphism
\[
(3)\qquad \varphi:\mathbb G_a\longrightarrow G
\]
\margin{4}

% Pied de page imprimé de l'archive.
\begin{center}
\small
4\\
\texttt{https://grothendieck.umontpellier.fr}\\
Université de Montpellier, UMR CNRS 5149, Case courrier 051, Place Eugène Bataillon\\
34095 Montpellier cedex 5 - France
\end{center}

\scanpage{6}

such that we have
\[
(4)\qquad \varphi(1)=g.
\]

I wonder whether this is also necessary, in ``good'' cases. Let us first take the
case where $S=\operatorname{Spec}K$. In characteristic $0$, we have a uniqueness
result for $\varphi$, which allows us to apply descent theory
(valid over any base \uline{in characteristic $0$}): if after an extension
$S'\longrightarrow S$ faithfully flat \unsure{quasi-compact}, there exists $\varphi'$ having
the desired properties, then $\varphi$ exists
(over $S$ itself $\unsure{\overline{\mathrm m}|=}$).
So over a \uline{field} of characteristic $0$ it is OK. In characteristic $p>0$ the difficulty
comes from the fact that there are group homomorphisms
\inserted{$\mathbb G_a\longrightarrow\mathbb G_a:$}
(such as $x\mapsto x^p-x$) which are \deleted{zero} nontrivial, and which are
trivial on $1$. We shall get around this \inserted{perhaps} by applying
\unsure{descent} with $\varphi$ which are \uline{monomorphisms}.
\deleted{Every} \inserted{\unsure{This raises the question:|This leads to the question:}}
Is a subgroup of $G$ \unsure{isomorphic} to $\mathbb G_a$ known when
we know an everywhere nontrivial section? Is this at least true over a
\uline{field} as base?

I do not feel like launching into these questions (certainly interesting!),
and prefer
\margin{5}

% Pied de page imprimé de l'archive.
\begin{center}
\small
5\\
\texttt{https://grothendieck.umontpellier.fr}\\
Université de Montpellier, UMR CNRS 5149, Case courrier 051, Place Eugène Bataillon\\
34095 Montpellier cedex 5 - France
\end{center}

\scanpage{7}
\mspage{651}

to limit myself to a case closer to the one I am after, $S=\operatorname{Spec}\mathbb Z$,
assuming $S$ noetherian integral, \deleted{\unsure{of characteristic}} with generic point $\eta$
of characteristic $0$. This already implies, by the result over a field of characteristic $0$,
that there exists a nonempty open $U$
\[
(5)\qquad
\underset{\substack{\text{open}\\\text{nonempty}}}{U\subset S}
\quad\text{and}\quad
\deleted{g_U}\;\varphi_U:\mathbb G_a\vert U\longrightarrow G\vert U
\]
such that we have (4) \uline{on} $U$ -- and (for $U$ chosen) this $\varphi_U$
is unique. By uniqueness, we thus see that there is a largest open
\inserted{$\unsure{\mathcal U|U}$} over which we can give such a
$\unsure{\varphi|\psi}_{\unsure{\mathcal U|U}}$,
and the question comes down to seeing whether we necessarily have $\unsure{\mathcal U|U}=S$.

If we had $U\ne S$, localizing if necessary at a maximal point of $\unsure{S\setminus\mathcal U|S\setminus U}$,
we \deleted{\unsure{can}} reduce to examining the case where $S$ is local,
with closed point $s$ -- and to \inserted{\unsure{trying to|a little of}} prove
that in this case, we can extend $\varphi_U$ over all of $S$
(whence a contradiction\ldots). In the case which interests us, $S$ is
\uline{regular} of dimension $1$, so by localizing we find a \uline{trait}.
I shall confine myself to this case, \unsure{crucial}
\margin{6}

% Pied de page imprimé de l'archive.
\begin{center}
\small
6\\
\texttt{https://grothendieck.umontpellier.fr}\\
Université de Montpellier, UMR CNRS 5149, Case courrier 051, Place Eugène Bataillon\\
34095 Montpellier cedex 5 - France
\end{center}

\scanpage{8}

\unsure{,|.} probably if we want to elucidate the general case, and which suffices for
my purpose.

Let $\eta$ be the generic point of $S$, $s$ its closed point. Consider the
subgroup $L_\eta=\varphi_\eta(\mathbb G_{a\eta})\subset G_\eta$,
let $L$ be its schematic closure in $G$. which is a reduced subgroup scheme
of $G$. Let \deleted{\illegible{1} $L$}
$L^0=L\setminus
\bigl(\underset{\substack{\text{components}\\\text{of }L_s\text{ non-}\\\text{identity}}}{\bigcup L_{s,i}}\bigr)$
be its identity component. We have \unsure{won} if we \unsure{prove} that
$L^0\simeq\mathbb G_{aS}$. In fact, we are not
$\unsure{\overline{\mathrm m}|=}$ sure that
$(L_s^0$ \deleted{\unsure{is}} $(=(L_s)^0)$ is reduced!
But for $\varphi:\mathbb G_{aS}\longrightarrow G$
\deleted{such that} extending $\varphi_\eta$ to exist, it is probably not necessary
that we have $L^0\simeq\mathbb G_a$ -- for it is possible that
$\varphi$ is not a monomorphism, i.e. that $\varphi_s$ is not a monomorphism --
it could $\unsure{\overline{\mathrm m}|=}$ be the trivial homomorphism
(if $g_s=1$)! But even if $g_s\ne1$, it is not clear that
$\varphi_s$ is a monomorphism.

Let us take the example where
\[
(6)\qquad G=\underline{\operatorname{Aut}}_{\mathcal O_S}(M)
\simeq\operatorname{Gl}(n)_S
\]
$M$ a \deleted{\unsure{a}} free module of rank $n$ over $S$,
then we shall have (since $\eta$ of characteristic $0$)
\[
(7)\qquad
\varphi_\eta(\lambda)=\sum\frac{X^{i}_{\unsure{\eta}}}{i!}\lambda^i
\]
where
\[
\unsure{X_\eta|X}\in\operatorname{End}_{\unsure{K}}(M_K)
\qquad K=k(\eta)
\]
\margin{7}

% Pied de page imprimé de l'archive.
\begin{center}
\small
7\\
\texttt{https://grothendieck.umontpellier.fr}\\
Université de Montpellier, UMR CNRS 5149, Case courrier 051, Place Eugène Bataillon\\
34095 Montpellier cedex 5 - France
\end{center}